\documentclass[../main.tex]{subfiles}
\begin{document}
\section{Pre-earnings premium}\label{sec:pre}

\subsection{Economic intuition}
Stocks have historically earned higher returns in the days around scheduled earnings announcements. Two explanations are compensation for the risk of the coming news, and buying pressure from investors whose attention the announcement draws \citep{frazzini2007,barber2013}. The strategy holds stocks into the announcement and exits just before the news, so it takes on the pre-announcement drift but not the announcement jump itself.

\subsection{Signal and trading rules}
The signal is the earnings calendar. Let $X$ be the last trading session before the announcement: the announcement day itself for an after-close release, and the day before for a before-open release. We buy the stock at the close five sessions before $X$ and sell at the close of $X$. Each position is 10\% of capital, hedged with an offsetting SPY position sized by the same 60-day beta as in Section~\ref{sec:drift}. We assume announcement dates are known at least seven trading days in advance, as companies normally publish them.

\subsection{Implementation}
The data cover 165 announcements: 105 in 2022--2023 and 60 in 2024. All trades are at the official close. The engine is \texttt{wave3/cpp/s\_events.cpp}.

\subsection{Results}
In 2022--2023 the hedged return averaged 0.46\% per event ($t = 1.47$), and the strategy earned 2.4\% a year over cash (Sharpe 0.88). In 2024 the average fell to 0.06\% per event ($t = 0.02$), with 0.4\% a year (Sharpe 0.17). Table~\ref{tab:pre} summarizes the results, and Figure~\ref{fig:events} compares the per-event averages with the drift strategy.

\begin{table}[ht]\centering\small
\caption{Pre-earnings premium, five-session hold, beta-hedged, after costs.}\label{tab:pre}
\begin{tabular}{lrrrrrr}
\toprule
 & \multicolumn{3}{c}{Per event} & \multicolumn{3}{c}{Account (annualized)} \\
\cmidrule(lr){2-4}\cmidrule(lr){5-7}
Period & Events & Mean & $t$-stat & Return & Sharpe & Max drawdown \\
\midrule
2022--2023 (fitting) & 105 & 0.46\% & 1.47 & 2.4\% & 0.88 & $-3.3\%$ \\
2024 (out-of-sample) & 60 & 0.06\% & 0.02 & 0.4\% & 0.17 & $-1.4\%$ \\
\bottomrule
\end{tabular}
\end{table}

Without the SPY hedge, the 2022--2023 average rises to 1.10\% per event ($t = 2.28$). Most of that is market exposure, and the unhedged 2024 average is close to zero.

\subsection{Assessment}
This is the weakest of the three. The 2022--2023 result is small and mostly disappears once its five best days are removed, and the premium did not appear in 2024. We keep it in the portfolio for its low volatility and low correlation with the other two strategies. Its own evidence is weak.

\ifSubfilesClassLoaded{% Bibliography.
\bigskip\noindent\emph{The reference details below should be checked against the published versions before citing.}\par\medskip
\begin{thebibliography}{99}
\setlength{\itemsep}{2pt}
\small
\bibitem[Bailey and L{\'o}pez de Prado(2014)]{bailey2014} Bailey, D. H., \& L{\'o}pez de Prado, M. (2014). The deflated Sharpe ratio. \emph{Journal of Portfolio Management}, 40(5).
\bibitem[Barber et~al.(2013)]{barber2013} Barber, B. M., De George, E. T., Lehavy, R., \& Trueman, B. (2013). The earnings announcement premium around the globe. \emph{Journal of Financial Economics}, 108(1).
\bibitem[Brandt et~al.(2008)]{brandt2008} Brandt, M. W., Kishore, R., Santa-Clara, P., \& Venkatachalam, M. (2008). Earnings announcements are full of surprises. Working paper.
\bibitem[Chan et~al.(1996)]{chan1996} Chan, L. K. C., Jegadeesh, N., \& Lakonishok, J. (1996). Momentum strategies. \emph{Journal of Finance}, 51(5).
\bibitem[Cremers and Weinbaum(2010)]{cremers2010} Cremers, M., \& Weinbaum, D. (2010). Deviations from put-call parity and stock return predictability. \emph{Journal of Financial and Quantitative Analysis}, 45(2).
\bibitem[DeMiguel et~al.(2009)]{demiguel2009} DeMiguel, V., Garlappi, L., \& Uppal, R. (2009). Optimal versus naive diversification: How inefficient is the 1/N portfolio strategy? \emph{Review of Financial Studies}, 22(5).
\bibitem[Easley et~al.(1998)]{easley1998} Easley, D., O'Hara, M., \& Srinivas, P. S. (1998). Option volume and stock prices: Evidence on where informed traders trade. \emph{Journal of Finance}, 53(2).
\bibitem[Fama and MacBeth(1973)]{famamacbeth1973} Fama, E. F., \& MacBeth, J. D. (1973). Risk, return, and equilibrium: Empirical tests. \emph{Journal of Political Economy}, 81(3).
\bibitem[Frazzini and Lamont(2007)]{frazzini2007} Frazzini, A., \& Lamont, O. A. (2007). The earnings announcement premium and trading volume. NBER Working Paper 13090.
\bibitem[Johnson and So(2012)]{johnsonso2012} Johnson, T. L., \& So, E. C. (2012). The option to stock volume ratio and future returns. \emph{Journal of Financial Economics}, 106(2).
\bibitem[Jorion(1986)]{jorion1986} Jorion, P. (1986). Bayes-Stein estimation for portfolio analysis. \emph{Journal of Financial and Quantitative Analysis}, 21(3).
\bibitem[Ledoit and Wolf(2004)]{ledoit2004} Ledoit, O., \& Wolf, M. (2004). A well-conditioned estimator for large-dimensional covariance matrices. \emph{Journal of Multivariate Analysis}, 88(2).
\bibitem[Markowitz(1952)]{markowitz1952} Markowitz, H. (1952). Portfolio selection. \emph{Journal of Finance}, 7(1).
\bibitem[Newey and West(1987)]{neweywest1987} Newey, W. K., \& West, K. D. (1987). A simple, positive semi-definite, heteroskedasticity and autocorrelation consistent covariance matrix. \emph{Econometrica}, 55(3).
\bibitem[Pan and Poteshman(2006)]{pan2006} Pan, J., \& Poteshman, A. M. (2006). The information in option volume for future stock prices. \emph{Review of Financial Studies}, 19(3).
\bibitem[Xing et~al.(2010)]{xing2010} Xing, Y., Zhang, X., \& Zhao, R. (2010). What does the individual option volatility smirk tell us about future equity returns? \emph{Journal of Financial and Quantitative Analysis}, 45(3).
\end{thebibliography}
}{}
\end{document}
